
\subsection{\texorpdfstring{\textbf{
      Реализация классической модели Хопфилда}}
  {Реализация классической модели Хопфилда}}

В связи с тем, что требуется быстро просматривать результаты
отдельно взятых операций, было решено использовать Jupyter Notebook.
При этом нужна была хорошая статическая типизация и скорость работы.
Scala 3 неплохо подходит под данные цели. Причём, у этого
языка программирования есть возможность использования библиотек
на Java, что, несомненно, является очень весомым преимуществом.

Был использован Almond Kernel \cite{almond_opt}
для запуска кода на Scala 3 в Jupyter Notebook. Также был создан докер
файл для независимого запуска кода с использованием
eclipse-temurin:21-jre-jammy. Была использована пара библиотек: Apache
Commons Math 3 \cite{apache_commons_math}
для работы с математикой и JFreeChart \cite{jfreechart}
для вывода графиков.

Весь код с записями можно найти по ссылке:

\url{https://github.com/paranormalKrono/intro-ml/blob/main/notebooks/Hopfield_Network.ipynb}

Далее, нужно закодировать картинки, преобразовывать изображения
в разные вектора с данными. Для Хэбба в вектор с бинарными данными,
так как стандартное правило Хебба поддерживает только бинарные вектора.

Были созданы разные кодировщики:
\begin{itemize}
  \item EncoderColorBits - переводит цвет в 24 элемента \(\{-1.0,1.0\}\)
  \item EncoderColorMono - переводит цвет в 1 элемент \(\{-1.0,1.0\}\)
\end{itemize}

Если будет 24 бита на цвет пикселя, то мы не теряем
никакой информации, но тогда получим вектор из 24 * 8 * 8
элементов = 1536. Стоит учитывать и тот факт, что сохраняющая способность
классической сети Хопфилда равна \(0,14N\). Получаем сохраняющую способность
равную 215,04 паттернов, но только в некореллированном случае. Иначе,
произойдёт значительное снижение ёмкости.

Получим закодированные изображения (Рис. \ref{media:images_patterns})
и изображения с ошибкой (Рис. \ref{media:images_patterns_error}).

\begin{figure}[h]
  \centering
  \includegraphics[width=0.5\textwidth]{image_patterns.png}
  \caption{}
  \label{media:images_patterns}
\end{figure}

\begin{figure}[h]
  \centering
  \includegraphics[width=0.5\textwidth]{image_patterns_error.png}
  \caption{}
  \label{media:images_patterns_error}
\end{figure}

Далее создана модель классической сети Хопфилда с дискретной
асинхронной и синхронной динамиками с правилом обучения по Хэббу.

Была проверена стационарность синхронной и асинхронной динамик -
паттерны, на которых была обучена модель не должны меняться.
Случай изначального расположения энергии в минимуме.

Для данных картинок с ошибками на входе асинхронная динамика уже через
2 * 24 * 8 * 8 (5 * 1536) шагов даёт изначальную картинку (Рис.
\ref{media:image_patterns_async}). Потребовалось больше шагов, так
как асинхронный шаг обновляет только один нейрон за шаг и для
корректного сравнения с синхронным методом были добавлены шаги.

\begin{figure}[h]
  \centering
  \includegraphics[width=0.4\textwidth]{image_patterns_async.png}
  \caption{}
  \label{media:image_patterns_async}
\end{figure}

Аналогично, синхронный метод даёт тот же результат за 1 шаг.

Далее было предложено проверить на цифрах
(Рис. \ref{media:image_patterns_numbers}) - монохромная кодировка.
И с монохромной ошибкой (Рис. \ref{media:image_patterns_numbers_error}).
Увы, результаты оказались значительно хуже. Возможно, из-за уменьшения
количества нейронов в 24 раза.

\begin{figure}[h]
  \centering
  \includegraphics[width=0.4\textwidth]{image_patterns_numbers.png}
  \caption{}
  \label{media:image_patterns_numbers}
\end{figure}

\begin{figure}[h]
  \centering
  \includegraphics[width=0.4\textwidth]{image_patterns_numbers_error.png}
  \caption{}
  \label{media:image_patterns_numbers_error}
\end{figure}

Также была попытка использовать другие кодировки,
но результаты были аналогичные.
